\chapter{Δεδομένα και Προεπεξεργασία}
\label{ch:data}

\section{Πηγές Δεδομένων}
\label{sec:data_sources}

Όλα τα δεδομένα αντλήθηκαν από το \textbf{ENTSO-E Transparency Platform}~\cite{entsoe_transparency},
το οποίο παρέχει ανοικτά δεδομένα για ευρωπαϊκές αγορές ηλεκτρισμού.
Συγκεκριμένα, χρησιμοποιήθηκαν οι παρακάτω πηγές:

\begin{itemize}
    \item \textbf{Day-Ahead Market τιμές (€/MWh):} Ωριαίες τιμές εκκαθάρισης της ελληνικής ζώνης
    υποβολής προσφορών (GR), για την περίοδο Ιανουαρίου 2017 -- Μαρτίου 2026.

    \item \textbf{Πραγματικό Ηλεκτρικό Φορτίο (MW):} Ωριαίο φορτίο ζήτησης για την Ελλάδα.

    \item \textbf{Day-Ahead Πρόβλεψη Φορτίου (MW):} Δημοσιευμένη από τον ΑΔΜΗΕ πριν την εκκαθάριση
    DAM~\cite{entsoe_load_dayahead}. Χρησιμοποιείται ως \textit{νόμιμο} μελλοντικό χαρακτηριστικό
    στην πρόβλεψη τιμής.

    \item \textbf{Παραγωγή ΑΠΕ (MW):} Πραγματική ωριαία παραγωγή ηλιακής, αιολικής, υδροηλεκτρικής,
    λιγνιτικής και αερίου παραγωγής, από τα δεδομένα Actual Generation per Production Type (A75).
    Τα δεδομένα παρέχονται σε 15-λεπτη ανάλυση και συναθροίζονται σε ωριαία.

    \item \textbf{Τιμή Φυσικού Αερίου TTF (€/MWh):} Ημερήσιες τιμές day-ahead της ολλανδικής αγοράς
    TTF, ως proxy του οριακού καυσίμου στην ελληνική αγορά. Για την περίοδο Νοεμβρίου 2025 --
    Φεβρουαρίου 2026, οι τιμές είναι συνθετικές (γραμμικά παρεμβαλλόμενες από anchor points).

    \item \textbf{Τιμή CO\textsubscript{2} (€/τόνο):} Ημερήσιες τιμές δικαιωμάτων εκπομπών EU ETS.
\end{itemize}

Το ολοκληρωμένο dataset περιλαμβάνει \textbf{80.126 ωριαίες παρατηρήσεις} (Ιαν.\ 2017 -- Μαρ.\ 2026).


\section{Χρονοσειρά Τιμών DAM}
\label{sec:price_timeseries}

Η χρονοσειρά τιμών της ελληνικής αγοράς DAM εμφανίζει έντονες εποχιακές συνιστώσες:
\begin{itemize}
    \item \textbf{Ωριαία εποχιακότητα:} Οι ώρες αιχμής (19:00--22:00) έχουν σταθερά υψηλότερες
    τιμές σε σχέση με τις ώρες κοιλάδας (3:00--6:00).
    \item \textbf{Εβδομαδιαία εποχιακότητα:} Τα Σαββατοκύριακα εμφανίζουν χαμηλότερο φορτίο και
    επομένως χαμηλότερες τιμές.
    \item \textbf{Ετήσια εποχιακότητα:} Υψηλές τιμές το καλοκαίρι (κλιματισμός) και τον χειμώνα
    (θέρμανση), χαμηλές την άνοιξη.
\end{itemize}

Χαρακτηριστικά στατιστικά της περιόδου αξιολόγησης (Q1 2026, Δεκ.\ 2025 -- Φεβ.\ 2026):

\begin{table}[H]
\centering
\caption{Στατιστικά περιγραφικά τιμής DAM ανά μήνα (Q1 2026).}
\label{tab:price_stats}
\begin{tabular}{lccccc}
\toprule
\textbf{Μήνας} & \textbf{Μέση τιμή (€)} & \textbf{Τ.Α.\ (€)} & \textbf{Ελάχ.\ (€)} & \textbf{Μέγ.\ (€)} & \textbf{Naive-1 MAE} \\
\midrule
Δεκ.\ 2025 & 110,4 & 33,6 & 0,0 & 210,0 & 10,86 \\
Ιαν.\ 2026 & 109,4 & 41,7 & 0,0 & 290,0 & 12,26 \\
Φεβ.\ 2026 & 77,8  & 45,4 & 0,0 & 275,0 & 13,60 \\
\bottomrule
\end{tabular}
\end{table}

Αξιοσημείωτο είναι το \textbf{regime shift} του Φεβρουαρίου 2026: η μέση τιμή πέφτει από 110€ σε
78€ (--29\%), ενώ η μεταβλητότητα αυξάνεται σε τυπική απόκλιση 45,4 €/MWh. Αυτό οφείλεται σε
συνδυασμό αυξημένης ηλιακής παραγωγής (μεταβατική περίοδος προς άνοιξη) και μεταβολής στο
ενεργειακό μίγμα.


\section{Χρονοσειρά Ηλεκτρικού Φορτίου}
\label{sec:load_timeseries}

Το ηλεκτρικό φορτίο της Ελλάδας χαρακτηρίζεται από σχετικά σταθερά εποχιακά πρότυπα:

\begin{table}[H]
\centering
\caption{Στατιστικά περιγραφικά φορτίου ανά μήνα (Q1 2026).}
\label{tab:load_stats}
\begin{tabular}{lcccc}
\toprule
\textbf{Μήνας} & \textbf{Μέση τιμή (MW)} & \textbf{Τ.Α.\ (MW)} & \textbf{Ελάχ.\ (MW)} & \textbf{Μέγ.\ (MW)} \\
\midrule
Δεκ.\ 2025 & 5.740 & 910 & 3.800 & 8.100 \\
Ιαν.\ 2026 & 5.820 & 980 & 3.700 & 8.400 \\
Φεβ.\ 2026 & 5.630 & 870 & 3.600 & 7.900 \\
\bottomrule
\end{tabular}
\end{table}

Σε αντίθεση με την τιμή, το φορτίο παρουσιάζει πιο ομαλή κατανομή και χαμηλότερη σχετική
μεταβλητότητα (Τ.Α./Μέση $\approx$ 15--17\% για φορτίο έναντι 30--58\% για τιμή), γεγονός που
εξηγεί γιατί η πρόβλεψη φορτίου επιτυγχάνει γενικά χαμηλότερο sMAPE.


\section{Κατασκευή Χαρακτηριστικών (Feature Engineering)}
\label{sec:features}

\subsection{Χρονικές Υστερήσεις (Lag Features)}

Οι χρονικές υστερήσεις της μεταβλητής στόχου αποτελούν τα σημαντικότερα χαρακτηριστικά για
ωριαίες χρονοσειρές. Χρησιμοποιήθηκαν:

\begin{itemize}
    \item \textbf{Lag-1:} Η τιμή ή το φορτίο της προηγούμενης ώρας.
    \item \textbf{Lag-2:} Η τιμή ή το φορτίο 2 ώρες πριν.
    \item \textbf{Lag-24:} Η ίδια ώρα της προηγούμενης ημέρας (ισχυρή ωριαία εποχιακότητα).
    \item \textbf{Lag-48:} Η ίδια ώρα 2 ημέρες πριν.
    \item \textbf{Lag-168:} Η ίδια ώρα της προηγούμενης εβδομάδας (εβδομαδιαία εποχιακότητα).
\end{itemize}

Για τα MIMO μοντέλα με \textit{dense lags}, προστέθηκαν επίσης οι υστερήσεις lag-3 έως lag-23,
δηλαδή τιμές κάθε ώρας μέσα στην ίδια ημέρα.

\subsection{Κυλιόμενα Στατιστικά (Rolling Features)}

\begin{itemize}
    \item \textbf{roll\_24:} Κυλιόμενος μέσος 24 ωρών.
    \item \textbf{roll\_48, roll\_168:} Κυλιόμενοι μέσοι 48 και 168 ωρών.
\end{itemize}

\subsection{Ημερολογιακά Χαρακτηριστικά}

\begin{itemize}
    \item \textbf{hour:} Ώρα ημέρας (0--23), ως κυκλική κωδικοποίηση $(\sin, \cos)$.
    \item \textbf{day\_of\_week:} Ημέρα εβδομάδας (0--6), ως κυκλική κωδικοποίηση.
    \item \textbf{month:} Μήνας (1--12), ως κυκλική κωδικοποίηση.
    \item \textbf{is\_weekend:} Δυαδικό (1 = Σ/Κ, 0 = εργάσιμη).
    \item \textbf{is\_holiday:} Δυαδικό για ελληνικές αργίες.
\end{itemize}

\subsection{Εξωγενείς Μεταβλητές}

\begin{table}[H]
\centering
\caption{Εξωγενή χαρακτηριστικά και η αντίστοιχη διαθεσιμότητά τους πριν DAM gate closure.}
\label{tab:exogenous}
\begin{tabular}{llcl}
\toprule
\textbf{Χαρακτηριστικό} & \textbf{Περιγραφή} & \textbf{Εργασία} & \textbf{Νόμιμο;} \\
\midrule
\texttt{load\_fc}       & Day-ahead πρόβλεψη φορτίου (MW) & Τιμή & Ναι \\
\texttt{gas\_price(D)}  & Τιμή TTF φυσικού αερίου (€/MWh) & Τιμή & Οριακά \\
\texttt{co2\_price(D)}  & Τιμή δικαιωμάτων CO₂ (€/τόνο) & Τιμή & Οριακά \\
\texttt{gen\_solar\_lag1} & Lag-1 ηλιακής παραγωγής (MW) & Τιμή, Φορτίο & Ναι \\
\texttt{gen\_wind\_lag1}  & Lag-1 αιολικής παραγωγής (MW) & Τιμή, Φορτίο & Ναι \\
\texttt{gen\_gas}       & Παραγωγή αερίου (MW, lagged) & Τιμή & Ναι \\
\texttt{gen\_lignite}   & Παραγωγή λιγνίτη (MW, lagged) & Τιμή & Ναι \\
\texttt{gen\_hydro}     & Παραγωγή υδρο/ηλεκ.\ (MW, lagged) & Τιμή & Ναι \\
\bottomrule
\end{tabular}
\end{table}

\noindent\textbf{Σχόλιο διαρροής:} Τα \texttt{gas\_price(D)} και \texttt{co2\_price(D)} αποτελούν
τιμές ίδιας ημέρας (day-ahead futures). Τεχνικά δημοσιεύονται πριν την εκκαθάριση DAM και αποτελούν
standard χαρακτηριστικά στη βιβλιογραφία EPF~\cite{lago2021}. Για την εργασία φορτίου, η τιμή DAM
δεν χρησιμοποιείται (μελλοντική άγνωστη κατά τον DAM), ενώ η παραγωγή ΑΠΕ αξιοποιείται μόνο ως
υστερήσεις (lagged values).


\section{Χρονικός Διαχωρισμός}
\label{sec:train_test_split}

Ακολουθείται αυστηρός χρονικός διαχωρισμός εκπαίδευσης--ελέγχου για αποφυγή
\textit{data leakage}~\cite{kaufman2012leakage}:

\begin{table}[H]
\centering
\caption{Χρονικός διαχωρισμός dataset.}
\label{tab:split}
\begin{tabular}{lll}
\toprule
\textbf{Σετ} & \textbf{Περίοδος} & \textbf{Σχόλιο} \\
\midrule
Εκπαίδευση & Ιαν.\ 2017 -- Νοε.\ 2025 & $\approx$76.000 παρατηρήσεις \\
Δοκιμή 1W  & Δεκ.\ 1--7, 2025 (168h) & Βραχυπρόθεσμη αξιολόγηση \\
Δοκιμή 1M  & Δεκ.\ 1--31, 2025 (744h) & Μηνιαία αξιολόγηση \\
Δοκιμή Q4  & Οκτ.\ 1 -- Δεκ.\ 31, 2025 (2208h) & Τριμηνιαία (εκπ.\ έως Σεπτ.\ 2025) \\
Δοκιμή Q1  & Δεκ.\ 1, 2025 -- Φεβ.\ 28, 2026 (2160h) & Κύρια αξιολόγηση \\
\bottomrule
\end{tabular}
\end{table}

Για τη στρατηγική Monthly Retraining (MR), η εκπαίδευση ανανεώνεται στο τέλος κάθε μήνα:
\begin{itemize}
    \item Δεκ.\ 2025: εκπαίδευση έως Νοε.\ 2025, πρόβλεψη Δεκ.\ 2025
    \item Ιαν.\ 2026: εκπαίδευση έως Δεκ.\ 2025, πρόβλεψη Ιαν.\ 2026
    \item Φεβ.\ 2026: εκπαίδευση έως Ιαν.\ 2026, πρόβλεψη Φεβ.\ 2026
\end{itemize}

\subsection{Σύνολο Χαρακτηριστικών}

Το τελικό σύνολο χαρακτηριστικών για κάθε εργασία:

\begin{table}[H]
\centering
\caption{Αριθμός χαρακτηριστικών ανά εργασία και παραλλαγή.}
\label{tab:feature_counts}
\begin{tabular}{lcc}
\toprule
\textbf{Παραλλαγή} & \textbf{Τιμή} & \textbf{Φορτίο} \\
\midrule
Standard (βασικά lags) & 152 & 133 \\
Dense (+ lags 3--23) & 170 & 151 \\
MIMO Two-Stage (+ 24 load forecasts) & 176 & --- \\
MIMO Two-Stage Dense & 194 & --- \\
\bottomrule
\end{tabular}
\end{table}
